
\documentclass[9pt,a4paper]{extarticle}

\usepackage[utf8]{inputenc}
\usepackage[T5]{fontenc}
\usepackage[vietnamese]{babel}
\usepackage[a4paper,top=6mm,bottom=6mm,left=6mm,right=6mm,columnsep=4mm]{geometry}
\usepackage{multicol}
\usepackage{enumitem}
\usepackage{titlesec}
\usepackage{microtype}
\usepackage{xcolor}
\usepackage{lmodern}
\usepackage{parskip}

\pagestyle{empty}
\setlength{\parindent}{0pt}
\setlength{\parskip}{1pt}
\setlength{\columnseprule}{0pt}

\linespread{0.92}

\titleformat{\section}{\bfseries\fontsize{9.2}{10}\selectfont}{}{0pt}{}
\titleformat{\subsection}{\bfseries\fontsize{8.6}{9.2}\selectfont}{}{0pt}{}
\titlespacing*{\section}{0pt}{2pt}{1pt}
\titlespacing*{\subsection}{0pt}{1pt}{0.5pt}

\setlist[itemize]{leftmargin=*,topsep=1pt,itemsep=0.5pt,parsep=0pt,partopsep=0pt}
\setlist[enumerate]{leftmargin=*,topsep=1pt,itemsep=0.5pt,parsep=0pt,partopsep=0pt}

\begin{document}
\small
\begin{center}
{\bfseries\fontsize{11}{12}\selectfont }\\[-1pt]
\end{center}

\vspace{-3pt}
\begin{multicols*}{2}

\section*{Câu 1. Sự ra đời, phát triển và vai trò của triết học Mác -- Lênin}
\subsection*{Nội dung rút gọn}
\begin{itemize}
  \item Triết học Mác -- Lênin ra đời vào khoảng thế kỷ XIX, không xuất hiện ngẫu nhiên mà do điều kiện lịch sử, xã hội và khoa học đã chín muồi.
  \item Chủ nghĩa tư bản phát triển mạnh làm sản xuất tăng lên, nhưng cũng làm xuất hiện mâu thuẫn gay gắt giữa giai cấp tư sản và giai cấp vô sản.
  \item Phong trào công nhân ngày càng lớn mạnh nên cần một hệ thống lý luận đúng đắn để chỉ đường.
  \item Về mặt tư tưởng, học thuyết này kế thừa có chọn lọc:
  \begin{itemize}
    \item triết học cổ điển Đức,
    \item kinh tế chính trị học Anh,
    \item chủ nghĩa xã hội không tưởng Pháp,
    \item thành tựu khoa học tự nhiên.
  \end{itemize}
  \item Mác và Ăngghen xây dựng nền tảng ban đầu; Lênin phát triển thêm trong hoàn cảnh mới, nhất là khi chủ nghĩa tư bản chuyển sang giai đoạn đế quốc chủ nghĩa.
  \item Đây là bước ngoặt cách mạng vì triết học không chỉ giải thích thế giới mà còn hướng con người đến việc cải tạo thế giới bằng thực tiễn.
  \item Hiện nay, triết học Mác -- Lênin vẫn có vai trò như một thế giới quan và phương pháp tư duy, giúp nhìn sự vật trong mối liên hệ, trong vận động và trong phát triển.
\end{itemize}

\textbf{Nhớ nhanh:} Ra đời từ thực tiễn xã hội + kế thừa tinh hoa cũ + hướng tới cải tạo thế giới.

\section*{Câu 2. Nguyên tắc phát triển; Mác vận dụng vào xã hội; sự phát triển của AI}
\subsection*{Nội dung rút gọn}
\begin{itemize}
  \item Mọi sự vật, hiện tượng đều luôn vận động, biến đổi, không đứng yên.
  \item Phát triển là sự thay đổi theo hướng từ thấp đến cao, từ đơn giản đến phức tạp, từ chưa hoàn thiện đến hoàn thiện hơn.
  \item Cơ sở lý luận nằm ở phép biện chứng duy vật: nguyên nhân sâu xa của phát triển nằm ngay trong mâu thuẫn bên trong sự vật.
  \item Chính sự đấu tranh giữa các mặt đối lập làm cho sự vật vận động và phát triển.
  \item Về phương pháp luận, cần:
  \begin{itemize}
    \item nhìn sự vật trong quá trình vận động,
    \item thấy xu hướng biến đổi,
    \item phân biệt cái mới và cái cũ,
    \item ủng hộ cái mới tiến bộ,
    \item tránh cách nhìn bảo thủ, đứng im.
  \end{itemize}
  \item Mác vận dụng nguyên tắc này vào nghiên cứu xã hội: xã hội không bất biến mà phát triển qua nhiều hình thái kinh tế -- xã hội.
  \item Động lực sâu xa của biến đổi xã hội là mâu thuẫn giữa lực lượng sản xuất và quan hệ sản xuất.
  \item AI cũng phát triển theo lịch sử: ý tưởng mô phỏng tư duy người $\rightarrow$ hệ chuyên gia $\rightarrow$ học máy $\rightarrow$ học sâu $\rightarrow$ AI tạo sinh.
  \item Sự phát triển của AI gắn với dữ liệu, thuật toán và năng lực tính toán ngày càng mạnh.
\end{itemize}

\textbf{Nhớ nhanh:} Muốn hiểu đúng một sự vật thì phải nhìn nó đang biến đổi như thế nào.

\section*{Câu 3. Mâu thuẫn lôgích và mâu thuẫn biện chứng; mâu thuẫn trong AI}
\subsection*{Nội dung rút gọn}
\begin{itemize}
  \item Mâu thuẫn lôgích và mâu thuẫn biện chứng là hai khái niệm khác nhau.
  \item \textbf{Mâu thuẫn lôgích} là lỗi trong tư duy hoặc phát biểu; ví dụ trong cùng điều kiện vừa khẳng định vừa phủ định cùng một điều.
  \item Loại mâu thuẫn này cần tránh vì làm lập luận thiếu chặt chẽ và mất sức thuyết phục.
  \item \textbf{Mâu thuẫn biện chứng} là mâu thuẫn có thật trong bản thân sự vật, là sự tồn tại đồng thời của các mặt đối lập vừa thống nhất vừa đấu tranh với nhau.
  \item Chính sự đấu tranh đó làm cho sự vật vận động, biến đổi và phát triển.
  \item Phép biện chứng là học thuyết về mâu thuẫn vì nó nghiên cứu sự vật trong trạng thái vận động, mà nguyên nhân sâu xa của vận động chính là mâu thuẫn bên trong.
  \item Không có mâu thuẫn thì không có động lực phát triển.
  \item Trong AI hiện nay có các mâu thuẫn nổi bật:
  \begin{itemize}
    \item phát triển nhanh nhưng phải an toàn,
    \item hiệu quả lớn nhưng có rủi ro đạo đức,
    \item cần dữ liệu lớn nhưng phải bảo vệ quyền riêng tư,
    \item tự động hóa cao nhưng ảnh hưởng đến việc làm và kỹ năng con người.
  \end{itemize}
  \item Vì vậy AI phát triển không theo đường thẳng đơn giản mà luôn gắn với tranh luận, điều chỉnh và hoàn thiện dần.
\end{itemize}

\textbf{Nhớ nhanh:} Mâu thuẫn lôgích là sai trong nghĩ; mâu thuẫn biện chứng là động lực phát triển của sự vật.

\section*{Câu 4. Cái mới, cái cũ và bài học trong CNTT}
\subsection*{Nội dung rút gọn}
\begin{itemize}
  \item \textbf{Cái mới} là cái ra đời phù hợp với xu hướng phát triển, có khả năng thay thế cái cũ và mở ra hướng tiến bộ hơn.
  \item \textbf{Cái cũ} là cái từng phù hợp trong một giai đoạn nhưng dần trở nên lạc hậu, thậm chí cản trở phát triển.
  \item Cuộc đấu tranh giữa cái mới và cái cũ luôn khó khăn, phức tạp, kéo dài vì cái cũ thường giữ vị trí quen thuộc, có lợi ích gắn với nó và được nhiều người bảo vệ.
  \item Trong khi đó, cái mới lúc đầu thường chưa hoàn thiện, chưa mạnh, chưa được số đông chấp nhận.
  \item Tuy vậy, nếu cái mới thực sự phù hợp với quy luật khách quan và đáp ứng nhu cầu phát triển thì về lâu dài nó sẽ thắng cái cũ.
  \item Trong CNTT:
  \begin{itemize}
    \item điện toán đám mây ban đầu bị nghi ngờ nhưng nay thành xu hướng lớn,
    \item AI trước đây bị xem là khó ứng dụng rộng rãi nhưng nay đã thâm nhập nhiều lĩnh vực,
    \item nhiều phương pháp phát triển phần mềm mới ban đầu không được tin tưởng nhưng sau đó được chấp nhận rộng rãi.
  \end{itemize}
  \item Bài học:
  \begin{itemize}
    \item biết phát hiện và ủng hộ cái mới tiến bộ,
    \item không bảo thủ bám mãi vào cái cũ,
    \item không nóng vội vì cái mới cần thời gian hoàn thiện,
    \item đánh giá cái mới bằng hiệu quả thực tế chứ không bằng cảm tính.
  \end{itemize}
\end{itemize}

\textbf{Nhớ nhanh:} Cái mới ban đầu thường yếu, nhưng nếu hợp xu thế thì cuối cùng vẫn thắng.

\section*{Câu 5. Ý thức phản ánh và sáng tạo thế giới; vai trò hiện nay; AI}
\subsection*{Nội dung rút gọn}
\begin{itemize}
  \item Ý thức phản ánh thế giới nghĩa là con người thông qua nhận thức có thể hiểu được hiện thực khách quan.
  \item Những gì con người biết, học, suy nghĩ và ghi nhận đều là sự phản ánh thế giới vào trong đầu óc mình.
  \item Nhưng ý thức không chỉ dừng ở phản ánh. Nó còn có khả năng sáng tạo vì con người biết đặt mục tiêu, lựa chọn cách làm, dự đoán kết quả và tổ chức hành động.
  \item Nhờ đó con người có thể biến những ý tưởng chưa có sẵn trong tự nhiên thành hiện thực thông qua lao động và thực tiễn.
  \item Trước khi có một cây cầu, phần mềm, ngôi nhà hay thiết bị công nghệ, chúng đều xuất hiện trước tiên trong tư duy dưới dạng mục tiêu, bản vẽ hay kế hoạch.
  \item Sau đó, bằng hoạt động thực tiễn, con người biến chúng thành sự vật có thật; theo nghĩa đó, ý thức góp phần sáng tạo ra thế giới của con người.
  \item Trong thời đại tri thức, khoa học và công nghệ, vai trò của ý thức càng lớn vì nhiều lĩnh vực phát triển phụ thuộc mạnh vào năng lực tư duy, sáng tạo và đổi mới.
  \item AI là thành tựu lớn của trí tuệ con người, đem lại nhiều lợi ích như tăng năng suất, hỗ trợ học tập, y tế, nghiên cứu và quản lý.
  \item Tuy nhiên, AI cũng đặt ra nhiều vấn đề: phụ thuộc công nghệ, thay đổi việc làm, nguy cơ lạm dụng, thông tin sai lệch, thiên vị thuật toán.
  \item Vì vậy cần nhìn AI một cách cân bằng: phát huy mặt tích cực nhưng phải kiểm soát rủi ro.
\end{itemize}

\textbf{Nhớ nhanh:} Ý thức không chỉ hiểu thế giới mà còn giúp con người biến đổi thế giới.

\section*{Câu 6. Nhận thức, chân lý và quá trình tiến bộ của nhân loại}
\subsection*{Nội dung rút gọn}
\begin{itemize}
  \item \textbf{Nhận thức} là quá trình con người phản ánh thế giới khách quan vào trong óc mình để hiểu thế giới.
  \item Quá trình này đi từ chưa biết đến biết, từ biết ít đến biết nhiều, từ hiểu chưa đúng đến hiểu đúng hơn.
  \item \textbf{Chân lý} là tri thức đúng đắn, phù hợp với hiện thực khách quan.
  \item Chân lý không phụ thuộc vào cảm xúc hay mong muốn chủ quan của con người.
  \item Tuy nhiên, chân lý con người đạt được ở mỗi giai đoạn thường có tính tương đối, đúng trong những điều kiện và phạm vi nhất định.
  \item Qua thời gian, nhận thức của con người ngày càng tiến gần hơn đến chân lý đầy đủ hơn.
  \item Có thể nói nhận thức của nhân loại về cơ bản là quá trình tích lũy ngày càng nhiều chân lý.
  \item Trong quá trình đó, con người cũng liên tục phát hiện sai lầm và sửa sai.
  \item Nhưng phát hiện sai lầm không phải mục đích cuối cùng; sai lầm chỉ là bước trung gian để đi đến cái đúng hơn.
  \item Ví dụ:
  \begin{itemize}
    \item khoa học từng cho rằng Mặt Trời quay quanh Trái Đất, sau này sửa lại,
    \item y học từng có nhiều quan niệm sai về bệnh tật, sau này chính xác hơn,
    \item trong công nghệ, nhiều mô hình cũ được thay bằng mô hình mới hiệu quả hơn.
  \end{itemize}
  \item Lịch sử nhận thức là quá trình vừa sửa sai vừa tiến bộ, trong đó chân lý được tích lũy ngày càng nhiều.
\end{itemize}

\textbf{Nhớ nhanh:} Con người tiến bộ bằng cách vừa tích lũy cái đúng, vừa sửa những cái sai.

\section*{Câu 7. Chân lý và thực tiễn theo quan điểm của Lênin}
\subsection*{Nội dung rút gọn}
\begin{itemize}
  \item Việc tư duy con người có đạt tới chân lý khách quan hay không không thể chỉ giải quyết bằng tranh luận lý thuyết.
  \item Muốn biết một nhận định có đúng hay không thì phải kiểm tra nó trong thực tiễn.
  \item Điều này thể hiện nội dung rất quan trọng: \textbf{thực tiễn là tiêu chuẩn của chân lý}.
  \item Thực tiễn ở đây bao gồm:
  \begin{itemize}
    \item hoạt động sản xuất vật chất,
    \item hoạt động chính trị -- xã hội,
    \item thực nghiệm khoa học,
    \item các hoạt động cải biến thế giới nói chung.
  \end{itemize}
  \item Nếu một lý thuyết khi đưa vào thực tế cho kết quả đúng, giúp giải quyết vấn đề và phù hợp hiện thực thì có cơ sở để xem nó là đúng.
  \item Ngược lại, nếu một lý thuyết nghe có vẻ hay nhưng không đứng vững trong thực tế thì cần xem xét lại.
  \item Câu nói này cũng cho thấy lý luận và thực tiễn phải thống nhất với nhau.
  \item Học để hiểu là cần, nhưng hiểu xong phải biết áp dụng, kiểm nghiệm và điều chỉnh.
  \item Quan điểm đó đồng thời phê phán lối tư duy chỉ suy đoán, xa rời đời sống thực tế.
\end{itemize}

\textbf{Nhớ nhanh:} Muốn biết đúng sai thật sự thì phải đem ra thực tế kiểm nghiệm.

\section*{Câu 8. Lý luận trở thành lực lượng vật chất khi thâm nhập vào quần chúng}
\subsection*{Nội dung rút gọn}
\begin{itemize}
  \item Câu nói này cho thấy mối quan hệ rất sâu sắc giữa lý luận và sức mạnh xã hội.
  \item Chỉ dùng lời nói hay sự phê phán trên giấy tờ thì chưa đủ để thay đổi hiện thực.
  \item Muốn thay đổi xã hội cần có lực lượng vật chất, tức là hành động thực tế của con người.
  \item Nhưng Mác cũng khẳng định rằng lý luận đúng có sức mạnh rất lớn.
  \item Khi một tư tưởng khoa học, tiến bộ và đúng đắn thâm nhập vào quần chúng, được tiếp nhận và biến thành hành động, thì lý luận không còn chỉ là ý tưởng trên sách vở mà trở thành sức mạnh vật chất thật sự.
  \item Điều đó cho thấy:
  \begin{itemize}
    \item vai trò to lớn của lý luận đúng đắn,
    \item vai trò quyết định của quần chúng nhân dân trong lịch sử,
    \item mối liên hệ giữa tư tưởng và hành động thực tiễn.
  \end{itemize}
  \item Nói cách khác, lý luận chỉ thực sự có sức mạnh khi đi vào đời sống và được con người thực hiện.
  \item Vì vậy, muốn biến tư tưởng thành hiện thực thì không chỉ cần lý luận đúng mà còn cần giáo dục, tổ chức và vận động quần chúng.
\end{itemize}

\textbf{Nhớ nhanh:} Lý luận đúng + quần chúng hành động = sức mạnh vật chất.

\section*{Câu 9. Thống nhất lý luận và thực tiễn; vận dụng trong CNTT}
\subsection*{Nội dung rút gọn}
\begin{itemize}
  \item Lý luận và thực tiễn có mối quan hệ chặt chẽ với nhau, không thể tách rời.
  \item \textbf{Lý luận} là sự khái quát kinh nghiệm thực tiễn thành tri thức có hệ thống.
  \item \textbf{Thực tiễn} là hoạt động vật chất có mục đích của con người nhằm cải biến tự nhiên và xã hội.
  \item Thực tiễn là cơ sở, nguồn gốc của lý luận vì mọi tri thức cuối cùng đều bắt đầu từ đời sống thực tế.
  \item Đồng thời, thực tiễn cũng là mục đích và là nơi kiểm tra tính đúng sai của lý luận.
  \item Ngược lại, lý luận có vai trò soi đường, định hướng và giúp hoạt động thực tiễn hiệu quả hơn.
  \item Yêu cầu phương pháp luận:
  \begin{itemize}
    \item xuất phát từ thực tiễn,
    \item tôn trọng thực tiễn,
    \item học đi đôi với hành,
    \item tránh lý thuyết suông,
    \item tránh làm theo kinh nghiệm mù quáng mà không có định hướng.
  \end{itemize}
  \item Trong CNTT:
  \begin{itemize}
    \item học lập trình không thể chỉ đọc lý thuyết mà phải viết code,
    \item thiết kế phần mềm phải xuất phát từ nhu cầu thực tế của người dùng,
    \item sản phẩm công nghệ phải được kiểm thử, sửa lỗi và cải tiến trong quá trình sử dụng,
    \item kiến thức trên lớp cần gắn với dự án, bài tập, sản phẩm cụ thể.
  \end{itemize}
  \item Nhờ thống nhất giữa lý luận và thực tiễn, việc học và làm trong CNTT mới đạt hiệu quả cao.
\end{itemize}

\textbf{Nhớ nhanh:} Lý luận phải đi với làm; học phải gắn với thực tế.

\section*{Câu 10. Xã hội phát triển qua các hình thái kinh tế -- xã hội}
\subsection*{Nội dung rút gọn}
\begin{itemize}
  \item Theo chủ nghĩa duy vật lịch sử, xã hội loài người không phát triển ngẫu nhiên mà theo những quy luật khách quan.
  \item Sự phát triển đó diễn ra thông qua các hình thái kinh tế -- xã hội khác nhau như cộng sản nguyên thủy, chiếm hữu nô lệ, phong kiến, tư bản chủ nghĩa và cao hơn nữa.
  \item Gọi đây là quá trình \textbf{lịch sử -- tự nhiên} vì:
  \begin{itemize}
    \item nó diễn ra trong lịch sử thật của loài người,
    \item đồng thời tuân theo quy luật khách quan, không phụ thuộc hoàn toàn vào ý muốn chủ quan của cá nhân nào.
  \end{itemize}
  \item Trong mỗi hình thái kinh tế -- xã hội, phương thức sản xuất giữ vai trò rất quan trọng.
  \item Khi lực lượng sản xuất phát triển đến một mức độ nhất định, quan hệ sản xuất cũ có thể không còn phù hợp nữa.
  \item Mâu thuẫn đó sẽ thúc đẩy xã hội biến đổi sang hình thái mới.
  \item Học thuyết hình thái kinh tế -- xã hội được xem là ``hòn đá tảng'' của chủ nghĩa duy vật lịch sử vì nó giúp giải thích xã hội một cách khoa học.
  \item Nó chỉ ra rằng muốn hiểu xã hội phải đi từ cơ sở vật chất, từ phương thức sản xuất, chứ không thể chỉ giải thích bằng ý chí cá nhân hay tư tưởng thuần túy.
  \item Nhờ học thuyết này, lịch sử xã hội được xem như một quá trình có quy luật, có nguyên nhân và có logic phát triển.
\end{itemize}

\textbf{Nhớ nhanh:} Xã hội phát triển theo quy luật khách quan, không phải ngẫu nhiên hay do ý muốn riêng lẻ.

\section*{Câu 11. Vì sao Ăngghen nói Mác tìm ra quy luật phát triển của xã hội?}
\subsection*{Nội dung rút gọn}
\begin{itemize}
  \item Darwin nổi tiếng vì đã phát hiện ra quy luật phát triển của thế giới sinh vật, cho thấy sinh giới không đứng yên mà tiến hóa qua thời gian.
  \item Ăngghen so sánh Mác với Darwin vì Mác cũng phát hiện ra quy luật phát triển của xã hội loài người.
  \item Trước Mác, nhiều người nghiên cứu xã hội nhưng thường nhìn lịch sử như một chuỗi sự kiện rời rạc hoặc giải thích bằng ý chí của các cá nhân lớn.
  \item Mác chỉ ra rằng xã hội phát triển theo những quy luật khách quan.
  \item Trong đó, phương thức sản xuất vật chất giữ vai trò nền tảng.
  \item Ông cũng chỉ ra mối quan hệ giữa:
  \begin{itemize}
    \item lực lượng sản xuất,
    \item quan hệ sản xuất,
    \item cơ sở hạ tầng,
    \item kiến trúc thượng tầng,
    \item đấu tranh giai cấp.
  \end{itemize}
  \item Chính những yếu tố đó làm cho xã hội vận động và biến đổi từ hình thái này sang hình thái khác.
  \item Đóng góp của Mác có ý nghĩa rất lớn: ông biến việc nghiên cứu xã hội từ chỗ chủ yếu mang tính cảm tính hoặc suy đoán thành một cách tiếp cận khoa học hơn.
\end{itemize}

\textbf{Nhớ nhanh:} Darwin tìm quy luật của sinh giới; Mác tìm quy luật của lịch sử xã hội.

\section*{Câu 12. Tư duy lý luận và sự thay đổi của triết học theo khoa học}
\subsection*{Nội dung rút gọn}
\begin{itemize}
  \item Ăngghen cho rằng một dân tộc muốn đứng trên đỉnh cao của khoa học thì không thể thiếu tư duy lý luận.
  \item Muốn phát triển khoa học không thể chỉ dựa vào kinh nghiệm rời rạc, mà cần khả năng khái quát, phân tích bản chất, phát hiện quy luật và liên hệ các tri thức với nhau.
  \item Tư duy lý luận giúp con người:
  \begin{itemize}
    \item nhìn xa hơn hiện tượng bề ngoài,
    \item hiểu bản chất sự vật,
    \item dự báo xu hướng phát triển,
    \item xây dựng các hệ thống tri thức lớn.
  \end{itemize}
  \item Nếu chỉ có kinh nghiệm thực tế mà thiếu tư duy lý luận, khoa học sẽ khó phát triển sâu và bền vững.
  \item Ăngghen cũng cho rằng mỗi lần khoa học đạt thành tựu mới thì triết học phải thay đổi hình thức tồn tại của mình.
  \item Điều đó nghĩa là triết học không thể đứng yên khi khoa học thay đổi.
  \item Những phát minh mới của khoa học buộc triết học phải điều chỉnh cách đặt vấn đề, cách giải thích và cách tồn tại của mình.
  \item Ví dụ, sự phát triển của vật lý, sinh học, công nghệ thông tin hay AI đều đặt ra những câu hỏi mới cho triết học về con người, nhận thức, tự nhiên và xã hội.
  \item Vì vậy, triết học chân chính phải luôn gắn bó với thành tựu khoa học, vừa khái quát chúng, vừa tự đổi mới để không lạc hậu.
\end{itemize}

\textbf{Nhớ nhanh:} Khoa học càng phát triển càng cần tư duy lý luận; khoa học đổi thì triết học cũng phải đổi.

\end{multicols*}
\end{document}
